\section{Síntesis y ampliación}

Las nueve sesiones de CS329A forman una cadena completa y progresiva: parten de los fundamentos de los LLM y los agents, mejoran la resolución de problemas en tiempo de prueba mediante muestreo, búsqueda, herramientas y verifiers; luego convierten las trayectorias exitosas en señales de entrenamiento como STaR, GRPO y DAPO; y después abordan la planeación a largo plazo, el deep research, la evaluación de trabajo real y los entornos abiertos.

La última sesión plantea cuatro preguntas de investigación que determinarán el avance de la siguiente etapa:

\begin{enumerate}
    \item \textbf{¿Cómo producir continuamente nuevo reasoning?} La multi-agent specialization mantiene la synthetic data diversity mediante diferencias estructuradas;
    \item \textbf{¿Cómo confiar en la retroalimentación?} DeepSeekMath-V2 usa meta-verification para auditar el análisis de problemas del verifier;
    \item \textbf{¿Cómo superar el banco de problemas?} Absolute Zero hace que el proposer y el solver generen conjuntamente un curriculum dinámico;
    \item \textbf{¿Cómo costear más cómputo?} Intelligence per watt integra en un mismo marco de optimización la capacidad, el hardware, el consumo energético y el serving local/cloud.
\end{enumerate}

El juicio de ingeniería que más vale conservar es este: la autosuperación no es que el modelo «se vuelva más inteligente por sí mismo» en el vacío, sino un ciclo cerrado restringido conjuntamente por el entorno, los datos, la búsqueda, la verificación, el entrenamiento, el sistema y la energía. El límite de capacidad del ciclo cerrado lo determinan a la vez el verifier más débil, la distribución de datos más estrecha y el recurso más costoso.

Por lo tanto, un agent de autosuperación de alta calidad debe contar con:

\begin{itemize}
    \item propuestas diversas y una cobertura de exploración medible;
    \item verifiers estratificados, independientes y refutables;
    \item un curriculum de tareas dinámico que varíe con la capacidad;
    \item un modelado de la incertidumbre para la retroalimentación lenta y el subjective reward;
    \item un límite claro entre memory y parameter learning;
    \item un enrutamiento local/cloud workload-aware y monitoreo de eficiencia energética;
    \item una calibración externa continua sobre entornos reales y objetivos humanos.
\end{itemize}

El curso termina, pero este campo cambia con enorme rapidez. Los artículos quedarán obsoletos; los principios de sistemas seguirán siendo reutilizables: ampliar el espacio de búsqueda, mejorar la calidad de la retroalimentación, elegir la señal de entrenamiento adecuada y convertir cada unidad de cómputo en más inteligencia confiable.

\end{document}
